% --- START OF FILE 04_activity_intents.tex ---

\chapter{Activity, Intent e Ciclo di Vita}

L'\textbf{Activity} è uno dei mattoni fondamentali (\textit{building blocks}) delle app Android. Rappresenta una singola schermata con un'interfaccia utente. Funge da punto di ingresso (\textit{entry point}) per l'interazione dell'utente con l'app.

\section{Gli Intent}
Un \textbf{Intent} è un oggetto di messaggistica (\textit{message-passing framework}) che permette a componenti diversi di richiedere azioni ad altri componenti, sia all'interno della stessa app che verso app esterne.

Possono essere di due tipi principali:
\begin{itemize}
    \item \textbf{Intent Espliciti (Explicit Intents):} Specificano \textit{esattamente} la classe del componente da avviare (es. \texttt{ActivityB::class.java}). Si usano quasi sempre all'interno della propria app.
    \item \textbf{Intent Impliciti (Implicit Intents):} Non specificano il componente esatto, ma indicano un'\textbf{azione} generica (es. \texttt{ACTION\_VIEW}) e dei dati (es. un URL). Il sistema Android verifica gli \textit{Intent-Filter} delle app installate e sceglie quella adatta (o chiede all'utente quale usare).
\end{itemize}

\begin{lstlisting}[caption={Esempi di Intent}, label={lst:intents}]
// 1. Intent Esplicito
val intentEsplicito = Intent(this, SecondActivity::class.java)
startActivity(intentEsplicito)

// 2. Intent Implicito (aprire una pagina web)
val uri = Uri.parse("http://www.uninsubria.it")
val intentImplicito = Intent(Intent.ACTION_VIEW, uri)
startActivity(intentImplicito)
\end{lstlisting}

\section{Activity Result API (Risultati dalle Activity)}
Spesso un'Activity ne avvia un'altra per ottenere un risultato (es. scattare una foto, scegliere un contatto).

\begin{tcolorbox}[colback=red!5!white,colframe=red!75!black,title=Attenzione all'esame: Deprecazione!]
    Il vecchio metodo \texttt{startActivityForResult()} e il relativo callback \texttt{onActivityResult()} sono \textbf{deprecati}. Erano problematici perché, se l'app veniva messa in background e uccisa dal sistema, il ritorno del risultato si perdeva.
\end{tcolorbox}

Il nuovo standard è l'\textbf{Activity Result API} (da \texttt{androidx 1.2.0}), che si divide in 3 passaggi fondamentali:
\begin{enumerate}
    \item \textbf{Contract:} Definisce input e output (es. \texttt{StartActivityForResult()}).
    \item \textbf{Register:} Si registra il contratto nella fase di inizializzazione chiamando \texttt{registerForActivityResult()}, che restituisce un \texttt{ActivityResultLauncher}.
    \item \textbf{Launch:} Si invoca il metodo \texttt{launch()} sull'oggetto Launcher quando si vuole avviare l'azione.
\end{enumerate}

\begin{lstlisting}[caption={Ottenere un risultato con Activity Result API (Esatto da Quiz)}]
// 1. Registrazione del contratto (da fare come attributo della classe)
val launcher = registerForActivityResult(
    ActivityResultContracts.StartActivityForResult()
) { result ->
    // Questo e il blocco dell'ActivityResultCallback
    if (result.resultCode == Activity.RESULT_OK) {
        val data = result.data?.getStringExtra("RISPOSTA")
    }
}

// 2. Lancio dell'Intent (es. dentro un listener)
button.setOnClickListener {
    val intent = Intent(this, SecondActivity::class.java)
    launcher.launch(intent) // L'istruzione per far partire la richiesta
}
\end{lstlisting}
\begin{tcolorbox}[colback=yellow!10!white,colframe=yellow!70!black,title=Perché \texttt{startActivityForResult} è deprecato?]
    Il vecchio metodo è stato deprecato perché il codice era strettamente accoppiato al ciclo di vita dell'Activity, generando bug (es. crash o perdita di dati) se l'Activity chiamante veniva distrutta dal sistema mentre la seconda Activity era aperta.
\end{tcolorbox}

\section{Il Ciclo di Vita dell'Activity (Lifecycle)}
Comprendere il ciclo di vita è essenziale. Android notifica i cambiamenti di stato dell'Activity invocando metodi di \textit{callback}.

\begin{center}
    \begin{tikzpicture}[
        node distance=0.8cm and 1cm,
        state/.style={rectangle, draw, rounded corners, fill=white, text width=3.5cm, align=center, drop shadow},
        arrow/.style={-Stealth, thick}
    ]
        \node[state, fill=blue!10] (create) {\textbf{onCreate()}\\Setup iniziale (UI)};
        \node[state, fill=cyan!10, below=of create] (start) {\textbf{onStart()}\\Activity Visibile};
        \node[state, fill=green!20, below=of start] (resume) {\textbf{onResume()}\\Focus / Interattiva};
        \node[state, fill=yellow!20, below=of resume] (pause) {\textbf{onPause()}\\Parz. visibile / No Focus};
        \node[state, fill=orange!20, below=of pause] (stop) {\textbf{onStop()}\\Completamente nascosta};
        \node[state, fill=red!20, below=of stop] (destroy) {\textbf{onDestroy()}\\Distruzione};

        \draw[arrow] (create) -- (start);
        \draw[arrow] (start) -- (resume);
        \draw[arrow] (resume) -- (pause);
        \draw[arrow] (pause) -- (stop);
        \draw[arrow] (stop) -- (destroy);

        \draw[arrow, dashed] (pause.east) -- ++(1.5,0) |- node[right, near start] {Ritorno in Foreground} (resume.east);
        \draw[arrow, dashed] (stop.west) -- ++(-1.5,0) |- node[left, near start] {Riavvio} (start.west);

        % Box laterali per le vite
        \draw[decorate, decoration={brace, amplitude=10pt}, thick]
        (resume.north east) -- (pause.south east) node[midway, right=15pt, align=left] {\textbf{Foreground Lifetime}\\(Multitasking, Popups)};

        \draw[decorate, decoration={brace, amplitude=10pt, mirror}, thick]
        (start.north west) -- (stop.south west) node[midway, left=15pt, align=right] {\textbf{Visible Lifetime}\\(L'utente vede l'app)};
    \end{tikzpicture}
\end{center}

Dettagli chiave sui metodi:
\begin{itemize}
    \item \textbf{\texttt{onCreate()}}: Chiamato una sola volta per instanziare l'UI (\texttt{setContentView}).
    \item \textbf{\texttt{onResume()} / \texttt{onPause()}}: Segnano i confini del \textit{Foreground lifetime}. In \texttt{onPause} l'app è ancora visibile ma non ha il focus (es. in modalità Multi-finestra, o se c'è un Dialog semi-trasparente sopra).
    \item \textbf{\texttt{onStop()}}: L'app non è più visibile (es. utente preme il tasto Home o apre un'altra app a schermo intero).
\end{itemize}

\section{Gestione delle Configuratrion Changes}
Quando avviene un cambiamento di configurazione a runtime (es. \textbf{rotazione dello schermo}, cambio della lingua, cambio modalità dark/light), il comportamento di default di Android è \textbf{distruggere e ricreare l'Activity} (\texttt{onPause} $\rightarrow$ \texttt{onStop} $\rightarrow$ \texttt{onDestroy} $\rightarrow$ \texttt{onCreate}).

\textbf{Il Problema:} Questo fa perdere tutte le variabili locali e lo stato non salvato.

Le soluzioni in ordine di correttezza (architetturale):
\begin{enumerate}
    \item \textbf{ViewModel (MVVM):} È l'approccio standard moderno. Il \texttt{ViewModel} "sopravvive" alle rotazioni. Lo stato della UI viene tenuto nel ViewModel.
    \item \textbf{\texttt{onSaveInstanceState()} e \texttt{onRestoreInstanceState()}:} Si salva un \texttt{Bundle} chiave-valore e lo si recupera alla ricreazione.
    \item \textbf{Forzare il Manifest (\texttt{android:configChanges}):} Specificando nel Manifest \texttt{android:configChanges="orientation|screenSize"}, Android \textit{non} riavvierà l'Activity, ma chiamerà \texttt{onConfigurationChanged()}. Questa pratica \textbf{non è consigliata} se non strettamente necessaria.
\end{enumerate}

\section{Tasks e Back Stack}
Un \textbf{Task} è un insieme di Activity con cui l'utente interagisce per svolgere un lavoro. Queste Activity sono organizzate in uno stack (LIFO) chiamato \textbf{Back Stack}.

Quando si avvia un'Activity con un Intent, si possono usare dei \textbf{Flags} per alterare il comportamento standard del Back Stack:
\begin{itemize}
    \item \texttt{Intent.FLAG\_ACTIVITY\_NEW\_TASK}: Avvia l'Activity in un nuovo task (o la porta in primo piano se il task esiste già). Indica che l'Activity funge da inizio di un nuovo lavoro.
    \item \texttt{Intent.FLAG\_ACTIVITY\_CLEAR\_TASK}: Causa la rimozione di tutte le Activity precedenti dal task corrente prima dell'avvio della nuova. Di solito usato in combinazione con \texttt{NEW\_TASK} (es. dopo un logout per impedire all'utente di tornare indietro).
\end{itemize}

\begin{center}
    \begin{tikzpicture}[
        node distance=0.5cm,
        box/.style={rectangle, draw, fill=blue!10, text width=2.5cm, align=center}
    ]
        \node[box] (a) {Activity A};
        \node[box, above=of a] (b) {Activity B};
        \node[box, above=of b, fill=green!20] (c) {Activity C (Top)};

        \draw[->, thick, red] (1.5, 1) -- node[right] {Tasto Back} (1.5, 0);
        \node[text width=5cm] at (5, 1) {Premendo "Back", Activity C viene distrutta (popped) e Activity B torna in stato di \texttt{onResume()}.};
    \end{tikzpicture}
\end{center}

\begin{tcolorbox}[colback=white,colframe=black,title=Comportamento dei Flag (Domande Quiz)]
    \begin{itemize}
        \item \textbf{\texttt{FLAG\_ACTIVITY\_NEW\_TASK}} e \textbf{\texttt{FLAG\_ACTIVITY\_CLEAR\_TASK}} usati insieme: Vengono rimosse tutte le Activity precedenti dal Task corrente e la nuova Activity diventa la root. Premendo "Back" si uscirà dall'app.
        \item Se l'utente preme il tasto Back: l'Activity in cima al Back Stack (foreground) viene distrutta (\texttt{onDestroy}) e la precedente torna in \texttt{onResume}.
    \end{itemize}
\end{tcolorbox}